\documentclass[]{article}

\usepackage{amsmath}
\usepackage{multirow}
\usepackage{natbib}
\usepackage{graphicx} 


\begin{document}

The determination of the neutrino mass hierarchy—whether Normal (NH) or Inverted (IH)—is a fundamental challenge in physics, with profound implications for cosmology and searches for neutrinoless double-beta decay. We address this question by computing the Bayesian evidence for each hierarchy, combining cosmological constraints on the sum of neutrino masses ($\Sigma m_\nu$) from the Dark Energy Spectroscopic Instrument (DESI) Data Release 2 with the latest neutrino oscillation data from NuFIT 6.0. To ensure the robustness of our conclusions against prior assumptions, we perform the analysis using two distinct frameworks: a physically-motivated hierarchical (SJPV) prior and an objective, information-theoretic (HS) reference prior. Within the standard $\Lambda$CDM cosmological model, the DESI DR2 data, which constrains $\Sigma m_\nu < 0.0642$ eV (95\% C.L.), places the minimum allowed mass for the IH ($\sim 0.099$ eV) in severe tension with observations. This results in decisive evidence for the Normal Hierarchy, with a Bayes factor ($K = P(D|\text{NH}) / P(D|\text{IH})$) of $K > 460$ even under the most conservative (HS) prior. We test the sensitivity of this conclusion to the cosmological model by extending the analysis to a $w_0w_a$CDM parameterization, finding that the preference, while reduced, remains strong ($K > 40$). The decisive preference for the NH implies a significantly more challenging landscape for upcoming neutrinoless double-beta decay experiments, as our posterior for the effective Majorana mass ($m_{\beta\beta}$) is suppressed into the few-meV range, well below the predictions for the now-disfavored Inverted Hierarchy.

\end{document}
                